\documentclass[11pt]{article}
\usepackage[margin=1in]{geometry}
\usepackage{amsmath,amssymb}
\usepackage{hyperref}
\usepackage[numbers,sort&compress]{natbib}
\usepackage{enumitem}

\title{Annotated Bibliography\\\large Set Theoretic Estimation and its Mechanization}
\author{Project \texttt{thejeb}}
\date{Compiled \today}

\begin{document}
\maketitle

\noindent
This bibliography tracks the sources underpinning the \texttt{thejeb}
mechanization of set theoretic estimation (STE). Entries are annotated
with (i) what the source establishes and (ii) its status for our Lean
formalization: \textbf{seed} (a paper we are directly mechanizing),
\textbf{foundational} (prior art we build on or cite), or
\textbf{survey/future} (candidate for later mechanization). Every claim
imported into the Lean development is cited back to a specific result
here; results are trusted only once machine-checked.

\section*{Seed papers}

\paragraph{\citet{combettes1993} --- \emph{The Foundations of Set
Theoretic Estimation}, Proc.\ IEEE 81(2):182--208.} \textbf{[seed]}
The synthesizing reference for STE. It reframes estimation away from
optimizing an objective function and toward \emph{feasibility}: a
solution is acceptable iff it is consistent with all available
information. Formally (\S II-C), each piece of information $i$ is a fuzzy
proposition $\Psi_i:\Xi\to[0,1]$ on the solution space $\Xi$; at a
confidence level $\psi_i$ it induces a \emph{property set}
$S_i=\{a\in\Xi\mid\Psi_i(a)\ge\psi_i\}$ (Eq.~3), and the
\emph{feasibility set} is $S=\bigcap_{i\in I}S_i$ (Eq.~4). A formulation
is \emph{consistent} if $S\ne\emptyset$, \emph{fair} if the true object
$h\in S$, and \emph{ideal} if $S=\{h\}$. \S III develops projection
algorithms (POCS) when $\Xi$ is a Hilbert space and the $S_i$ are closed
convex. \emph{Mechanization status:} the crisp core (property sets,
feasibility set, consistent/fair/ideal, monotonicity, invalid-information
detection) is machine-checked in \texttt{Ste.Basic}; the convex layer in
\texttt{Ste.Convex}. The fuzzy/graded case and the projection algorithms
are future work.

\paragraph{\citet{carlson2012} --- \emph{Set Theoretic Estimation
Applied to the Information Content of Ciphers and Decryption}, Ph.D.\
dissertation, University of Idaho.} \textbf{[seed]}
Recasts STE over a \emph{topological} solution space rather than the
Hilbert space of prior work, avoiding the Optimal Bounding Ellipsoid
(OBE) pathology whereby the bounding volume can spuriously expand the
solution set. Applies STE to decryption: the solution space is a finite
key space, each intercepted ciphertext induces a property set of keys
under which it decodes to a meaningful message, and decryption succeeds
when the feasible key set collapses (Shannon unicity). Named results we
extracted (see \texttt{papers/notes}): Lemma~2.1 (error bound on
spurious keys), Lemma~2.2 (property-set intersection no larger than the
smallest set), Theorem~2.3 (property sets are AEP typical sets, tying
information theory under STE), Lemma~4.1 (an $S$-cipher has
$(|A|-|T|)!$ isomorphic keys), Theorem~4.1 (key count for a $P$-cipher),
Theorem~5.1 (PS/SP product ciphers vs.\ block substitution unicity),
Theorem~5.2 (a permutation cipher reduces to a substitution cipher
within a symbol), Corollary~5.3 (a boundary-aligned PSP cipher is
idempotent to an $S$-cipher; hence all block ciphers are block
substitution ciphers). \emph{Mechanization status:} the STE-decryption
skeleton (cipher systems, key property sets, feasible keys, fairness of
genuine traffic, unicity as ideal information, observation monotonicity)
is machine-checked in \texttt{Ste.Cipher}; the counting theorems
(4.1, Lemma~4.1) are the most tractable next targets, the unicity/AEP
results the deepest.

\section*{Foundational}

\paragraph{\citet{shannon1949} --- \emph{Communication Theory of Secrecy
Systems}.} \textbf{[foundational]}
Origin of unicity distance and key equivocation, which Carlson recasts
in set theoretic terms. The set theoretic ``unicity'' we formalize
(feasible key set $=\{k_0\}$) is the qualitative core of Shannon's
quantitative unicity point.

\paragraph{\citet{youla1982} --- \emph{Image Restoration by the Method
of Convex Projections}.} \textbf{[foundational]}
POCS: iterated projection onto closed convex constraint sets converges
to a feasible point. The algorithmic engine behind Combettes~\S III and
the natural sequel to \texttt{Ste.Convex}.

\paragraph{\citet{bregman1967} --- relaxation / common point of convex
sets.} \textbf{[foundational]}
The cyclic-projection method predating POCS; the ancestor of convex
feasibility solvers.

\paragraph{\citet{schweppe1968} --- set-membership state estimation;
\citet{deller1989} --- set-membership identification in DSP.}
\textbf{[foundational]}
Early ``unknown but bounded'' set-membership estimation---one of the
scattered traditions Combettes unifies; the source of the OBE algorithms
Carlson critiques.

\section*{Survey / future mechanization}

\paragraph{\citet{bauschke1996} --- \emph{On Projection Algorithms for
Solving Convex Feasibility Problems}, SIAM Review.} \textbf{[survey]}
The definitive survey of the convex feasibility problem
$\text{find }x\in\bigcap_i C_i$. Its firmly-nonexpansive-operator
framework is the target for a machine-checked convergence proof of POCS
atop \texttt{Ste.Convex}; Mathlib already carries the requisite
inner-product-space projection theory.

\paragraph{\citet{combettes1993numerical} --- \emph{The Convex
Feasibility Problem in Image Recovery}.} \textbf{[survey]}
Companion to the seed paper; concrete constraint sets and numerical
behavior for the convex STE instantiation.

\section*{Tooling}

\paragraph{\citet{moura2021} --- Lean 4; \citet{mathlib2020} ---
Mathlib.} \textbf{[tooling]}
The proof assistant and library in which every result in this project is
machine-checked. We pin Mathlib \texttt{v4.32.0} (toolchain
\texttt{leanprover/lean4:v4.32.0}).

\bibliographystyle{plainnat}
\bibliography{refs}

\end{document}
